\documentclass[../../../include/open-logic-section]{subfiles}

\begin{document}

\olfileid[tr]{sth}{ordinals}{ordtype}
\olsection{Sıra Tipleri Olarak Sıra Sayıları}

Yerine Koyma'yla donanıp artık $\ZFminus$ içinde çalıştığımıza göre,
nihayet amaçladığımız sonucu ispatlayabiliriz:

\begin{thm}\ollabel{thmOrdinalRepresentation}
Her iyi sıralama biricik bir sıra sayısıyla izomorftur.
\end{thm}

\begin{proof}
$\tuple{A,<}$ bir iyi sıralama olsun. \olref[basic]{ordisoidentity}
gereği en çok bir sıra sayısıyla izomorftur. Öyleyse olmayana ergiyle
$\tuple{A,<}$'nin \emph{hiçbir} sıra sayısıyla izomorf olmadığını varsayın.
Önce ``$\tuple{A,<}$'yi olabildiğince küçülteceğiz''. Ayrıntılı olarak:
bir $\tuple{A_a,<_a}$ öz başlangıç kesiti hiçbir sıra sayısıyla izomorf
değilse bu özelliği taşıyan en küçük bir $a\in A$ vardır; o zaman
$B=A_a$ ve $\mathord{\lessdot}=\mathord{<_a}$ olsun. Böyle bir kesit yoksa
$B=A$ ve $\mathord{\lessdot}=\mathord{<}$ olsun.

Tanım gereği $B$'nin her öz başlangıç kesiti bir sıra sayısıyla izomorftur;
bu sıra sayısı yukarıdaki gibi biriciktir. Dolayısıyla Yerine Koyma gereği
şu küme vardır ve bir fonksiyondur:
\[
	f = \Setabs{\tuple{\beta, b}}{b \in B\text{ ve }\ordeq{\beta}{\tuple{B_b, \lessdot_b}}}
\]
Olmayana ergiyi tamamlamak için bir $\alpha$ sıra sayısı için $f$'nin
$\alpha\to B$ izomorfizmi olduğunu göstereceğiz.

$\ran{f}=B$ olduğu açıktır. \olref[iso]{lemordsegments} gereği $f$
sıralamayı korur; yani $\gamma\in\beta$ ancak ve ancak
$f(\gamma)\lessdot f(\beta)$ olur. $\dom{f}$'nin bir sıra sayısı olduğunu
göstermek için \olref[basic]{corordtransitiveord} gereği $\dom{f}$'nin
geçişli olduğunu göstermek yeterlidir. Bir $\beta\in\dom{f}$ sabitleyin;
yani bir $b$ için
$\ordeq{\beta}{\tuple{B_b,\lessdot_b}}$ olsun. $\gamma\in\beta$ ise
\olref[iso]{wellordinitialsegment} gereği $\gamma\in\dom{f}$ olur;
genelleyerek $\beta\subseteq\dom{f}$ elde edilir.
\end{proof}

Bu sonuç, \olref[vn]{sec} kesiminden beri sunmak istediğimiz şu tanıma izin
verir:

\begin{defn}
$\tuple{A,<}$ bir iyi sıralamaysa onun sıra tipi $\ordtype{A,<}$,
$\ordeq{\tuple{A,<}}{\alpha}$ koşulunu sağlayan biricik $\alpha$ sıra
sayısıdır.
\end{defn}

Ayrıca bu tanım iki güzel ilkeye izin verir:

\begin{cor}\ollabel{ordtypesworklikeyouwant}
$\tuple{A,<}$ ve $\tuple{B,\lessdot}$ iyi sıralamalarsa:
\begin{align*}
	\ordtype{A, <} = \ordtype{B, \lessdot}&\text{ ancak ve ancak }\ordeq{\tuple{A, <}}{\tuple{B, \lessdot}}\\
	\ordtype{A, <} \in \ordtype{B, \lessdot}&\text{ ancak ve ancak bir }b \in B\text{ için }\ordeq{\tuple{A, <}}{\tuple{B_b, \lessdot_b}}.
\end{align*}
\end{cor}

\begin{proof}
Özdeşlik \olref[basic]{ordisoidentity} gereği geçerlidir. İkinci iddiayı
ispatlamak için $\ordtype{A,<}=\alpha$ ve
$\ordtype{B,\lessdot}=\beta$ olsun; izomorfizmimiz
$f\colon\beta\to\tuple {B,\lessdot}$ olsun. O zaman:
\begin{align*}
	\alpha \in \beta&\text{ ancak ve ancak }\funrestrictionto{f}{\alpha} \colon \alpha \to B_{f(\alpha)}\text{ bir izomorfizmdir}\\
	&\text{ ancak ve ancak }\ordeq{\tuple{A, <}}{\tuple{B_{f(\alpha)}, \lessdot_{f(\alpha)}}}\\
	&\text{ ancak ve ancak bir $b \in B$ için }\ordeq{\tuple{A, <}}{\tuple{B_b, \lessdot_b}}
\end{align*}
olur; burada \olref[iso]{ordisounique},
\olref[iso]{wellordinitialsegment} ve
\olref[basic]{ordissetofsmallerord} kullanılmıştır.
\end{proof}

\end{document}
